% GENERATED FILE. Edit canonical lessons, then run npm run generate:books.
% canonical-source-hash: fnv1a64:7750175990a5070c
% canonical-lessons: PT-C32-costas, PT-C32-cabelo, PT-C32-barriga, PT-C32-rosto, PT-C32-sangue

\chapter{Back, Hair, Belly, Face, Blood}
\label{ch:pt-back-hair-belly-face-blood}
\hlchaptermodality{32}

\begin{chapteropening}
\textbf{By the end of this chapter:} \emph{I can name the back, hair, the belly, the face and blood.}

Complete the last lesson of chapter 32: I can name the back, hair, the belly, the face and blood.
\end{chapteropening}

\section[as costas]{as costas — the back}

\label{lesson:PT-C32-costas}



\begin{quote}
Before the new one: say the Portuguese for a shoulder, then the Portuguese for a knee.
\end{quote}

\subsection*{You'll want to know: as costas}
\textbf{as costas} — \textquotedblleft{}the back\textquotedblright{}. Plural in form, like English \textquotedblleft{}trousers\textquotedblright{}.

\begin{cousinweb}
\textbf{costas} $\leftarrow$ Latin \emph{costa}, \textquotedblleft{}rib, side\textquotedblright{} (EtymDB 2.1) — inherited. The plural came to name the whole back.
\end{cousinweb}

\begin{grammarlens}[title={What you've built}]
One more word for this chapter.
\end{grammarlens}

\subsection*{Your turn}
\begin{itemize}
\raggedright
  \item \emph{Say it:} \emph{as costas}
  \item \emph{Say it:} \emph{as costas}, once more
  \item \emph{Say it:} say \emph{as costas}
  \item \emph{Recall:} say the Portuguese for a shoulder, then the Portuguese for a knee, then say \emph{as costas} again
\end{itemize}

\begin{tcolorbox}[breakable,colback=teal!4,colframe=teal!35!black,title={Before you move on}]
What does as costas mean? (\textquotedblleft{}The back\textquotedblright{}.) Say it once more.
\end{tcolorbox}
\section[o cabelo]{o cabelo — hair (on the head)}

\label{lesson:PT-C32-cabelo}



\begin{quote}
Before the new one: say the Portuguese for a knee, then the Portuguese for the back.
\end{quote}

\subsection*{You'll want to know: o cabelo}
\textbf{o cabelo} — \textquotedblleft{}hair (on the head)\textquotedblright{}.

\begin{cousinweb}
\textbf{cabelo} $\leftarrow$ Latin \emph{capillus}, \textquotedblleft{}hair of the head\textquotedblright{} (EtymDB 2.1) — inherited. The Latin \emph{-p-} between vowels became \emph{-b-}.
\end{cousinweb}

\begin{grammarlens}[title={What you've built}]
One more word for this chapter.
\end{grammarlens}

\subsection*{Your turn}
\begin{itemize}
\raggedright
  \item \emph{Say it:} \emph{o cabelo}
  \item \emph{Say it:} \emph{o cabelo}, once more
  \item \emph{Say it:} say \emph{o cabelo}
  \item \emph{Recall:} say the Portuguese for a knee, then the Portuguese for the back, then say \emph{o cabelo} again
\end{itemize}

\begin{tcolorbox}[breakable,colback=teal!4,colframe=teal!35!black,title={Before you move on}]
What does o cabelo mean? (\textquotedblleft{}Hair (on the head)\textquotedblright{}.) Say it once more.
\end{tcolorbox}
\section[a barriga]{a barriga — the belly}

\label{lesson:PT-C32-barriga}



\begin{quote}
Before the new one: say the Portuguese for the back, then the Portuguese for hair.
\end{quote}

\subsection*{You'll want to know: a barriga}
\textbf{a barriga} — \textquotedblleft{}the belly\textquotedblright{}.

\begin{cousinweb}
\textbf{barriga} is probably built from Old Portuguese \emph{barrica}, \textquotedblleft{}cask\textquotedblright{} (EtymDB 2.1): a belly shaped like a little barrel. That link is the usual account, not a certain one.
\end{cousinweb}

\begin{grammarlens}[title={What you've built}]
One more word for this chapter.
\end{grammarlens}

\subsection*{Your turn}
\begin{itemize}
\raggedright
  \item \emph{Say it:} \emph{a barriga}
  \item \emph{Say it:} \emph{a barriga}, once more
  \item \emph{Say it:} say \emph{a barriga}
  \item \emph{Recall:} say the Portuguese for the back, then the Portuguese for hair, then say \emph{a barriga} again
\end{itemize}

\begin{tcolorbox}[breakable,colback=teal!4,colframe=teal!35!black,title={Before you move on}]
What does a barriga mean? (\textquotedblleft{}The belly\textquotedblright{}.) Say it once more.
\end{tcolorbox}
\section[o rosto]{o rosto — the face}

\label{lesson:PT-C32-rosto}



\begin{quote}
Before the new one: say the Portuguese for hair, then the Portuguese for the belly.
\end{quote}

\subsection*{You'll want to know: o rosto}
\textbf{o rosto} — \textquotedblleft{}the face\textquotedblright{}.

\begin{cousinweb}
\textbf{rosto} $\leftarrow$ Latin \emph{rostrum}, \textquotedblleft{}beak, snout\textquotedblright{} — inherited. EtymDB 2.1 has no usable row for \emph{rosto}; it derives Spanish \emph{rostro} from Latin \emph{rostrum}.
\end{cousinweb}

\begin{grammarlens}[title={What you've built}]
One more word for this chapter.
\end{grammarlens}

\subsection*{Your turn}
\begin{itemize}
\raggedright
  \item \emph{Say it:} \emph{o rosto}
  \item \emph{Say it:} \emph{o rosto}, once more
  \item \emph{Say it:} say \emph{o rosto}
  \item \emph{Recall:} say the Portuguese for hair, then the Portuguese for the belly, then say \emph{o rosto} again
\end{itemize}

\begin{tcolorbox}[breakable,colback=teal!4,colframe=teal!35!black,title={Before you move on}]
What does o rosto mean? (\textquotedblleft{}The face\textquotedblright{}.) Say it once more.
\end{tcolorbox}
\section[o sangue]{o sangue — blood}

\label{lesson:PT-C32-sangue}



\begin{quote}
Before the new one: say the Portuguese for the belly, then the Portuguese for the face.
\end{quote}

\subsection*{You'll want to know: o sangue}
\textbf{o sangue} — \textquotedblleft{}blood\textquotedblright{}.

\begin{cousinweb}
\textbf{sangue} $\leftarrow$ Latin \emph{sanguis}, \textquotedblleft{}blood\textquotedblright{}, through Vulgar Latin \emph{*sanguem} — inherited. EtymDB 2.1 has no ancestor row for \emph{sangue}; it derives Romanian \emph{sânge} from \emph{*sanguem}.
\end{cousinweb}

\begin{grammarlens}[title={What you've built}]
One more word for this chapter.
\end{grammarlens}

\subsection*{Your turn}
\begin{itemize}
\raggedright
  \item \emph{Say it:} \emph{o sangue}
  \item \emph{Say it:} \emph{o sangue}, once more
  \item \emph{Say it:} say \emph{o sangue}
  \item \emph{Recall:} say the Portuguese for the belly, then the Portuguese for the face, then say \emph{o sangue} again
\end{itemize}

\begin{tcolorbox}[breakable,colback=teal!4,colframe=teal!35!black,title={Before you move on}]
What does o sangue mean? (\textquotedblleft{}Blood\textquotedblright{}.) Say it once more.
\end{tcolorbox}
